\section{Synthetic Gaussian projections}\label{sec:synthetic}

The Gaussian test in Section~\ref{sec:gaussian} uses an orthonormal
selector of the projected difference space's complement.  Here we instead
draw additional Gaussian rows in that complement.  Their composition with
the side matrix has a random covariance, which must be averaged in the
likelihood comparison.  To normalize that comparison, we apply the
Gaussian-base equal-label identity to the projected prior, rather than
applying only its spherical instance to the original prior.

The inverse-determinant scale of this random covariance cancels the scale
of the projected simplex volume.  A critical inverse-distance kernel then
gives a double-logarithmic density cost.  Although the cost has the same
order as in Section~\ref{sec:gaussian}, the comparison law and cancellation
are different.  We use this bound while conditioning on an entire state
history.

We use the spherical and Gaussian notation of Section~\ref{sec:replicas}.
Fix \(d\ge16\) and put
\begin{equation}\label{syn:parameters}
 k=\ell=\left\lfloor\frac d8\right\rfloor,\qquad
 u=\ell-1,\qquad m=\left\lfloor\frac d2\right\rfloor,\qquad q=m-u.
\end{equation}
In particular,
\begin{equation}\label{syn:dimension-margins}
 k+u\le m-2,\qquad q-k-1=m-2k>0,\qquad
 d-1-u-k=d-2k\ge\frac{3d}{4},\qquad
 \frac{m}{q-k-1}\le2.
\end{equation}
The letter \(m\) in this section is the dimension of the side projection.
The difference count in Theorem~\ref{thm:equal-label-measure} will be \(u\).

\begin{proposition}[Information beyond an independent Gaussian projection]
\label{syn:block}
Let \(P=f\sigma_d\) be a probability, where \(f\) is a nonnegative Borel
function satisfying \(f\le\Lambda<\infty\), with \(\Lambda\ge1\).
Draw \(X\sim P\) and \(A\sim\gamma_{k,d}\) independently.  Let
\(z=z(A,AX)\) be a measurable function taking at most \(K\ge1\) values.
Independently of this entire experiment draw \(B\sim\gamma_{m,d}\), and put
\(Y=BX\).  With the parameters in \eqref{syn:parameters},
\begin{equation}\label{syn:block-bound}
 I(X;z\mid B,Y)
 \le \frac{\log K}{\ell}
       +C\bigl[d+\log(1+\log\Lambda)\bigr],
\end{equation}
where \(C\) is absolute and independent of \(P,\Lambda,K\), and the message
rule.
\end{proposition}

The proof couples \(\ell\) signals by giving them the same actual
measurements.  Corollary~\ref{cor:replica-likelihood} describes their joint
law.  We compare its total correlation before and after applying \(B\).
For the comparison after \(B\), a Gaussian matrix is constrained to
annihilate the projected differences.  We first construct the resulting
probability law, then estimate its density while keeping its dependence
on \(B\) explicit.

\subsection{The actual and projected replica laws}

For a tuple \(\mathbf x=(x_1,\ldots,x_\ell)\) in \(\R^h\), where
\(h\) will be \(d\) or \(m\), write
\[
 D_{\mathbf x}=[x_2-x_1,\ldots,x_\ell-x_1],\qquad
 L_{\mathbf x}=\col(D_{\mathbf x}),\qquad
 \Delta(\mathbf x)=\det(D_{\mathbf x}^{\mathsf T}D_{\mathbf x})^{1/2}.
\]
Thus \(L_{\mathbf x}\) and \(\Delta(\mathbf x)\) are the objects
\(F_{\mathbf x}\) and \(J(\mathbf x)\) of
\eqref{eq:equal-label-geometry}, with difference count \(u\).  At full
affine rank define
\begin{equation}\label{syn:tuple-weight}
 w(\mathbf x)=(2\pi)^{-ku/2}\Delta(\mathbf x)^{-k}.
\end{equation}

We need the same version of a projection density before and after a
composition of matrices.  If \(\nu\) is a finite measure and \(C\) has
\(r\) rows, write
\begin{equation}\label{syn:density-version}
 g_C^\nu(v)=\liminf_{j\to\infty}
       \int\varphi_{1/j,r}(Cs-v)\,\nu(ds),\qquad
 \varphi_{\eta,r}(v)=(2\pi\eta^2)^{-r/2}
               \exp\!\left(-\frac{\norm v^2}{2\eta^2}\right).
\end{equation}
This formula defines a jointly measurable function even for matrices of
deficient rank.  In every use below the matrix has full row rank almost
surely, and Theorem~\ref{thm:equal-label-measure} identifies this version
with the projection density and transfers its finite positive values to the
constrained evaluations.  We abbreviate \(g_C^P\) to \(g_C\).

Fix the data of Proposition~\ref{syn:block}.  In
Corollary~\ref{cor:replica-likelihood}, take ambient dimension \(d\),
\(k\) rows, difference count \(u\), and prior \(P\).  The condition
\(k+u\le d-2\) follows from \eqref{syn:dimension-margins}.
Draw \(A\sim\gamma_{k,d}\), then a common label \(b\) with density
\(g_A\), and finally draw
\(X_1,\ldots,X_\ell\) independently from \(P\) conditional on \(AX=b\),
and set \(z=z(A,b)\).  Denote this law by \(Q\).  Its tuple--matrix
likelihood is
\begin{equation}\label{syn:actual-replica}
 Q(d\mathbf x,dA)
 =P^{\otimes\ell}(d\mathbf x)\gamma_{k,d}^{L_{\mathbf x}}(dA)
       w(\mathbf x)g_A(Ax_1)^{-u}.
\end{equation}
Every pair \((X_i,z)\) has the original signal--message law.
Independently draw \(B\) and put \(Y_i=BX_i\).  For fixed \(B\), let
\(P^B=B_\#P\), and define
\[
 \mathcal C_X=\KL(Q_{\mathbf X}\Vert P^{\otimes\ell}),\qquad
 \mathcal C_Y=\E_B\KL(Q_{\mathbf Y\mid B}\Vert(P^B)^{\otimes\ell}).
\]
Lemma~\ref{lem:replica-logs}, with the same substitution and density
bound \(L=\Lambda\), proves absolute integrability of
\(\log g_A(b)\) and \(\log w(\mathbf X)\).  It gives
\begin{equation}\label{syn:correlation-upper}
 \mathcal C_X
 \le \E\log w(\mathbf X)-u\,\E\log g_A(b)<\infty.
\end{equation}
Data processing under the coordinatewise map \(x_i\mapsto Bx_i\) gives
\(0\le\mathcal C_Y\le\mathcal C_X\).

We next construct a probability density against which to test
\(\mathcal C_Y\).  Almost every \(B\) has full row rank.  For such a fixed
\(B\), the spherical projection formula
\eqref{eq:spherical-projection-density}, with \(m\) rows, and
\(P\le\Lambda\sigma_d\) show that \(P^B\) has bounded Lebesgue density
on a compact ellipsoid.  Here \(d-m-2\ge0\).  The standard Gaussian density
has a positive minimum on that ellipsoid, so \(P^B\) has bounded density
relative to standard Gaussian probability on \(\R^m\).  This bound is
finite for each such \(B\) and is allowed to depend on \(B\).

Apply Corollary~\ref{cor:replica-likelihood} in ambient dimension \(m\),
again with \(k\) rows and difference count \(u\).  Its dimension condition
is \(k+u\le m-2\), as in \eqref{syn:dimension-margins}.  The resulting
same-label tuple law has density
\begin{equation}\label{syn:comparison-density}
 R_B(\mathbf y)
 =w(\mathbf y)\,
   \E_{F\sim\gamma_{k,m}^{L_{\mathbf y}}}g_{FB}(Fy_1)^{-u}
\end{equation}
relative to \((P^B)^{\otimes\ell}\).  Indeed, the density version agrees
pointwise under composition:
\begin{equation}\label{syn:composition-version}
 g_F^{P^B}(v)
 =\liminf_j\int\varphi_{1/j,k}(Fy-v)\,P^B(dy)
 =g_{FB}^{P}(v).
\end{equation}
The finite-measure theorem transfers version, zero, and infinity
exceptions before the negative power is taken.  Thus \(R_B\) is finite and
positive almost everywhere for the product law and integrates to one.
It can be chosen jointly measurable in \(B,\mathbf y\): on the
full-affine-rank set, realize the projected Gaussian rows with the Borel
orthogonal projection onto \(L_{\mathbf y}^{\perp}\), and assign fixed
values on the null complement.

The actual law \(Q_{\mathbf Y\mid B}\) is absolutely continuous relative
to \((P^B)^{\otimes\ell}\), because \(Q_{\mathbf X}\ll P^{\otimes\ell}\)
and \(B\) is independent of the tuple.  At the actual \((B,\mathbf Y)\),
draw \(F\) with the conditional law in \eqref{syn:comparison-density}
and put
\[
 \Xi=g_{FB}(FY_1).
\]
The same null-set transfer makes \(\Xi\) finite and positive almost surely
in this enlarged experiment.

\subsection{The synthetic density and its critical kernel}

The comparison density \(R_B\) contains the inverse \(u\)th moment of
\(\Xi\).  To compare its logarithm with the actual likelihood in
\eqref{syn:correlation-upper}, we will control \(\Xi\) using the actual
matrix \(A\) and label \(b\).  The next lemma first averages the synthetic
Gaussian covariance at a fixed tuple.  It then bounds the resulting
inverse-distance integral by a function \(U_A(b)\) whose total mass can
be estimated independently of that tuple.

\begin{lemma}[Synthetic likelihood and critical kernel]\label{syn:synthetic}
Let \(P=f\sigma_d\) be as in Proposition~\ref{syn:block}.  Fix a tuple
\(\mathbf x\in(S^{d-1})^\ell\) of affine rank \(u\), and put
\[
 x=x_1,\qquad L=L_{\mathbf x},\qquad
 \delta(s)=\dist(s-x,L).
\]
Draw \(B\sim\gamma_{m,d}\).  Given \(B\), draw
\(F\sim\gamma_{k,m}^{BL}\), so each row of \(F\) is a standard Gaussian
in \((BL)^\perp\).  Then
\begin{equation}\label{syn:synthetic-bound}
 \E_{B,F}g_{FB}(FBx)
 \le (2\pi)^{-k/2}a_q\int\delta(s)^{-k}\,P(ds),\qquad
 a_q=\E\det(HH^{\mathsf T})^{-1/2}\le(q-k-1)^{-k/2},
\end{equation}
where \(H\) is a standard Gaussian \(k\times q\) matrix.

There is an absolute \(C_0>0\) such that the following holds with
\begin{equation}\label{syn:cutoff}
 r_*=\exp\!\left[-\frac{C_0(d+\log\Lambda)}{d-1-u-k}\right]\in(0,1).
\end{equation}
For every full-row-rank \(k\times d\) matrix \(A\), put \(a=\norm A_{\op}\)
and define
\begin{align}
 U_A(b)={}&(2\pi)^{-k/2}\ind_{\{\norm b\le a\}}
 \left[1+a^k\int g_A(v)\ind_{\{\norm{v-b}\le2a\}}\right. \notag\\[-2mm]
 &\hspace{42mm}\left.
       {}\cdot\min\{(ar_*)^{-k},\norm{v-b}^{-k}\}\,dv\right].
 \label{syn:majorant}
\end{align}
At \(v=b\), the second entry in the minimum is interpreted as \(+\infty\).
If \(Ax=b\) and \(L\subset\ker A\), then
\begin{equation}\label{syn:kernel-bound}
 (2\pi)^{-k/2}\int\delta(s)^{-k}\,P(ds)\le U_A(b).
\end{equation}
Writing \(V_k\) for the volume of the Euclidean unit ball in \(\R^k\),
the same function satisfies
\begin{align}
 \int U_A(b)\,db
 &\le (2\pi)^{-k/2}V_k a^k[2+k\log(2/r_*)],
 \label{syn:majorant-mass}\\
 \E_{A\sim\gamma_{k,d}}\log\int U_A(b)\,db
 &\le C\bigl[d+\log(1+\log\Lambda)\bigr].
 \label{syn:majorant-log}
\end{align}
The logarithm in the last display is absolutely integrable, and
\(\E_A\log(1+\int U_A)<\infty\).
\end{lemma}

\begin{proof}
At the fixed tuple, the restriction of \(B\) to \(L\) has rank \(u\)
almost surely.  Choose an orthonormal \(m\times q\) matrix \(N\) whose
columns span \((BL)^\perp\), using only \(BD_{\mathbf x}\).  Gram--Schmidt
with a fixed ordering gives a measurable choice on the full-rank set.
With an independent standard Gaussian \(k\times q\) matrix \(H\), the
required conditional law of \(F\) is realized by
\begin{equation}\label{syn:synthetic-realization}
 F=HN^{\mathsf T}.
\end{equation}
The choice of \(N\) uses only \(B|_L\).  The Gaussian restriction
\(B|_{L^\perp}\) is independent of \(B|_L\).  Write \(s-x=v+w\) with
\(v\in L\) and \(w\in L^\perp\).  Then \(N^{\mathsf T}Bv=0\), while,
conditionally on \(B|_L\), the vector \(N^{\mathsf T}Bw\) has law
\(N(0,\norm w^2I_q)\).  Therefore, conditionally on \(H\) and \(B|_L\),
\begin{equation}\label{syn:covariance}
 FB(s-x)\sim N(0,\delta(s)^2HH^{\mathsf T}).
\end{equation}
This conditional calculation accounts for the dependence of \(F\) on \(B\).

For \(\eta>0\), Tonelli's theorem and \eqref{syn:covariance} give
\[
 \E_{B,F}\int\varphi_{\eta,k}(FB(s-x))\,P(ds)
 =(2\pi)^{-k/2}\int
   \E_H\det(\delta(s)^2HH^{\mathsf T}+\eta^2I_k)^{-1/2}\,P(ds).
\]
The integrand on the right increases as \(\eta\downarrow0\).  The tube
bound \eqref{eq:affine-tube} gives \(\delta(s)>0\) for \(P\)-almost every
\(s\); it also gives a finite inverse \(k\)th moment, since
\(d-1-u>k\).  The matrix \(H\) has full row rank almost surely.
Monotone convergence on the right and Fatou's lemma on the left, along
\(\eta=1/j\), prove the first inequality in \eqref{syn:synthetic-bound}.

For completeness, successive orthogonal residuals of the rows of \(H\)
have independent squared lengths
\[
 V_i\sim\chi^2_{q-i+1},\qquad
 \det(HH^{\mathsf T})=\prod_{i=1}^kV_i.
\]
Indeed, conditionally on the preceding rows, the new row is standard
Gaussian on a complement of dimension \(q-i+1\); this conditional law is
constant, which also proves independence of the residual lengths.  This is
the elementary Gaussian orthogonal reduction described in
Edelman and Rao~\cite[Section~5]{EdelmanRao2005}; the rectangular product
used here follows from the displayed conditioning.  Direct integration of
the chi-squared density gives \(\E V^{-1}=1/(p-2)\) for
\(V\sim\chi_p^2\), \(p>2\).  Cauchy--Schwarz hence gives
\[
 a_q=\prod_{i=1}^k\E V_i^{-1/2}
 \le\prod_{i=1}^k(q-i-1)^{-1/2}
 \le(q-k-1)^{-k/2}.
\]

We next use the actual observations.  If \(Ax=b\) and \(L\subset\ker A\),
then, for \(v=P_L(s-x)\),
\begin{equation}\label{syn:actual-kernel}
 \norm{As-b}=\norm{A(s-x-v)}\le a\,\delta(s).
\end{equation}
Let \(\alpha=d-1-u\).  The tube bound, now applied to \(x+L\), gives
\[
 P\{\delta\le t\}\le\Lambda C^d t^\alpha,\qquad 0<t\le1.
\]
Layer-cake integration therefore yields, for \(0<r\le1\),
\[
 \int_{\{\delta\le r\}}\delta^{-k}\,dP
 \le \Lambda C^d
       \left(1+\frac{k}{\alpha-k}\right)r^{\alpha-k}.
\]
Since \(\alpha-k=d-2k\ge3d/4\), an absolute choice of \(C_0\) makes this
quantity at most one at \(r=r_*\).  On \(\{\delta>r_*\}\),
\eqref{syn:actual-kernel} gives both
\(\delta^{-k}\le r_*^{-k}\) and
\(\delta^{-k}\le a^k\norm{As-b}^{-k}\).
The pushforward of \(P\) under \(A\) has density \(g_A\), and both \(As\)
and \(b\) have norm at most \(a\).  Integration of these two bounds gives
\eqref{syn:kernel-bound} with the two support cutoffs in
\eqref{syn:majorant}.

The remaining singularity has the critical exponent \(k\) in dimension
\(k\).  Radial integration gives the exact mass
\begin{equation}\label{syn:critical-mass}
 \int_{\{\norm w\le2a\}}
       \min\{(ar_*)^{-k},\norm w^{-k}\}\,dw
 =V_k[1+k\log(2/r_*)].
\end{equation}
The ball of radius \(ar_*\) contributes \(V_k\), and the annulus contributes
\(kV_k\int_{ar_*}^{2a}dt/t\).  Dropping the outer cutoff in
\eqref{syn:majorant} only for its convolution term, and using
\(\int g_A=1\), proves \eqref{syn:majorant-mass}.

We record the scale in this bound before taking its logarithm.  Integrating
the standard Gaussian density over the ball of radius \(\sqrt k\) gives
\((2\pi)^{-k/2}V_k\le(e/k)^{k/2}\).  Also
\(\E\norm A_{\op}\le C\sqrt d\): choose \(1/4\)-nets on the two unit
spheres of cardinalities at most \(9^k\) and \(9^d\), bound the operator
norm by twice the maximum of the scalar bilinear forms on the nets, and
integrate the union bound for their standard Gaussian tails.  Jensen's
inequality gives
\[
 \E\log\bigl[(2\pi)^{-k/2}V_k\norm A_{\op}^k\bigr]
 \le\frac k2(1-\log k)+k\log(C\sqrt d)\le Cd.
\]
The cutoff \eqref{syn:cutoff} and \eqref{syn:dimension-margins} give
\[
 \log[2+k\log(2/r_*)]
 \le C+\log d+\log(1+\log\Lambda).
\]
These estimates prove \eqref{syn:majorant-log}.  Finally,
\(\log\norm A_{\op}\) is integrable below because the norm dominates the
absolute value of one standard Gaussian entry.  Retaining the first term
of \(U_A\) bounds its total mass below by
\((2\pi)^{-k/2}V_k\norm A_{\op}^k\).
Together with \eqref{syn:majorant-mass}, these observations prove the two
remaining integrability assertions.
\end{proof}

\subsection{Integrability and cancellation of the correlation costs}

The majorant now controls the synthetic density at the actual label.
Before subtracting the two correlation bounds, we verify integrability
of the comparison logarithms.  This uses both the finite actual
correlation and the normalization of \(R_B\).

\begin{lemma}[Integrability of the projected comparison]\label{syn:logs}
In the replica experiment \(Q\), enlarged by the independent \(B\) and
the conditional \(F\) above,
\[
 \E|\log w(\mathbf Y)|+\E|\log\Xi|
       +\E|\log R_B(\mathbf Y)|<\infty.
\]
If \(G\) is a standard Gaussian \(m\times u\) matrix, then
\begin{align}
 \E\log w(\mathbf Y)
 &=\E\log w(\mathbf X)-\frac{k}{2}\E\log\det(G^{\mathsf T}G),
 \label{syn:gram-log}\\
 \frac12\E\log\det(G^{\mathsf T}G)&\le\frac u2\log m.
 \label{syn:gram-upper}
\end{align}
Moreover,
\begin{equation}\label{syn:correlation-lower}
 \mathcal C_Y
 \ge\E\log R_B(\mathbf Y)
 \ge\E\log w(\mathbf Y)-u\,\E\log\Xi.
\end{equation}
\end{lemma}

\begin{proof}
At a fixed full-affine-rank \(\mathbf X\), write
\(D_{\mathbf X}=UR\), where \(U\) has orthonormal columns.  Since \(B\)
is independent, \(BU\) is a standard Gaussian \(m\times u\) matrix.
Consequently
\[
 \Delta(\mathbf Y)\stackrel{\mathrm{law}}{=}
 \Delta(\mathbf X)\det(G^{\mathsf T}G)^{1/2}
 \quad\text{conditionally on }\mathbf X.
\]
The logarithm of \(\Delta(\mathbf X)\) is integrable by
Lemma~\ref{lem:replica-logs}.  Gaussian column residuals factor
\(\det(G^{\mathsf T}G)\) into independent chi-squared variables of degrees
\(m,m-1,\ldots,m-u+1\).  Their absolute logarithms are integrable from
their densities at zero and infinity.  This proves
\eqref{syn:gram-log} and the asserted integrability of \(\log w(\mathbf Y)\).
Jensen applied to each factor gives \eqref{syn:gram-upper}.

The actual equal-label experiment has \(L_{\mathbf X}\subset\ker A\).
Conditionally on \((A,b,\mathbf X)\), the matrix \(B\) remains independent,
and the conditional law of \(F\) is
\(\gamma_{k,m}^{B L_{\mathbf X}}\).  Lemma~\ref{syn:synthetic} therefore
gives
\begin{equation}\label{syn:conditional-mean}
 \E[\Xi\mid A,b,\mathbf X]\le a_qU_A(b).
\end{equation}
It follows that \(\E\log(1+\Xi)<\infty\).  Here are the details needed for
the possible small values of \(g_A(b)\).  Conditional Jensen reduces the
claim to integrability of \(\log(1+U_A(b))\).  On the actual finite positive
set for \(g_A\),
\[
 \log(1+U_A(b))
 \le\log(1+g_A(b))+\log(1+U_A(b)/g_A(b)).
\]
The first term is integrable by Lemma~\ref{lem:replica-logs}.  For the second,
average over \(b\) with density \(g_A\) and apply Jensen to obtain at most
\(\log(1+\int U_A)\).  Lemma~\ref{syn:synthetic} makes its average finite.
Thus \((\log\Xi)^+\) is integrable.

Let the product reference first draw \(B\) and then draw \(\ell\)
independent points from \(P^B\).  The divergence of the actual joint law
from this reference is \(\mathcal C_Y<\infty\), and \(R_B\) integrates to one under this
reference.  Equation~\eqref{eq:entropy-positive-log} gives
\begin{equation}\label{syn:positive-comparison-log}
 \E(\log R_B(\mathbf Y))^+
 \le\E\log(1+R_B(\mathbf Y))\le\mathcal C_Y+\log2<\infty.
\end{equation}
Conditionally on \(B,\mathbf Y\), Equation
\eqref{syn:comparison-density} has a finite positive right side, so
\(\E_F\Xi^{-u}<\infty\).  Hence
\(\E_F(\log\Xi)^-<\infty\) on almost every such fiber; also
\(\E_F(\log\Xi)^+<\infty\) there by the preceding joint bound.
Conditional Jensen is now legitimate and gives
\[
 \log R_B(\mathbf Y)
 \ge\log w(\mathbf Y)-u\,\E_F\log\Xi.
\]
Rearranging this inequality,
\[
 u\,\E_F(\log\Xi)^-
 \le(\log R_B(\mathbf Y))^+
       +|\log w(\mathbf Y)|+u\,\E_F(\log\Xi)^+.
\]
Its right side has finite expectation, proving joint integrability of
\((\log\Xi)^-\).  The same Jensen lower bound gives
\[
 (\log R_B(\mathbf Y))^-
 \le|\log w(\mathbf Y)|+u\,\E_F|\log\Xi|,
\]
so the negative comparison logarithm is integrable as well.
Finally use the normalized test \(R_B\) in \eqref{eq:entropy-test}.
All logarithms are now integrable, and that test together with the
conditional Jensen bound proves \eqref{syn:correlation-lower}.
\end{proof}

\begin{proof}[Proof of Proposition~\ref{syn:block}]
The preceding lemma permits subtraction of all likelihood and correlation
terms.  Conditional Jensen in \eqref{syn:conditional-mean} gives
\[
 \E[\log\Xi\mid A,b,\mathbf X]\le\log a_q+\log U_A(b).
\]
The last logarithm is integrable under the actual label law: \(U_A(b)\)
is bounded below there by \((2\pi)^{-k/2}\), and its positive logarithm
was controlled in the proof of Lemma~\ref{syn:logs}.  At fixed \(A\),
Jensen with \(b\sim g_A\) gives
\[
 \int g_A(b)\log\frac{U_A(b)}{g_A(b)}\,db
 \le\log\int U_A(b)\,db.
\]
Together with \eqref{syn:synthetic-bound} and \eqref{syn:majorant-log},
this yields
\begin{equation}\label{syn:likelihood-difference}
 \E\log\Xi-\E\log g_A(b)
 \le-\frac k2\log(q-k-1)
       +C[d+\log(1+\log\Lambda)].
\end{equation}
Combine \eqref{syn:correlation-upper}, \eqref{syn:gram-log},
\eqref{syn:gram-upper}, \eqref{syn:correlation-lower}, and
\eqref{syn:likelihood-difference}.  The result is
\begin{align}
 \mathcal C_X-\mathcal C_Y
 &\le\frac{ku}{2}\log\frac{m}{q-k-1}
       +Cu[d+\log(1+\log\Lambda)] \notag\\
 &\le C\ell[d+\log(1+\log\Lambda)].
 \label{syn:correlation-loss}
\end{align}
The first term is \(O(k u)\) by \eqref{syn:dimension-margins}.
Thus the scale from the projected Gaussian determinant cancels the
inverse-determinant scale in the synthetic likelihood.  The critical
kernel was integrated before taking a logarithm, leaving only
\(\log(1+\log\Lambda)\).

Apply Lemma~\ref{lem:replica-information} with \(t=\ell\),
the representative \(S=X_1\), message \(W=z\), and map
\(\phi_B(x)=Bx\).  Every replica has the required signal--message pair law,
\(B\) is independent of the whole tuple and message, and
\(\mathcal C_X<\infty\).  Its conclusion and \(H(z)\le\log K\) give
\[
 \ell I(X_1;z\mid B,BX_1)
 \le\log K+\mathcal C_X-\mathcal C_Y.
\]
The joint law of \((B,X_1,z)\) is the product of the \(B\)-law and the
original \((X,z)\)-law.  Substitution of \eqref{syn:correlation-loss}
therefore proves \eqref{syn:block-bound}.
\end{proof}

\subsection{Iteration along the full state history}

\begin{proposition}[The synthetic route to the sample bound]\label{syn:streaming}
For every nonnegative integer-valued \(M(d)=o(d^2)\) and every sequence
\(0<\epsilon(d)\le1/10\), consider learners in
Definition~\ref{def:streaming-model} with deterministic horizon \(T(d)\)
and signal \(X\sim\sigma_d\).  If their angular
success probability, averaged over the uniform signal, rows, and learner
randomness, is at least \(2/3\), then
\[
 T(d)\ge c\,d\log(1/\epsilon(d))
\]
for all sufficiently large \(d\), with an absolute \(c>0\).  The threshold
may depend on the memory sequence.
\end{proposition}

\begin{proof}
Apply Lemma~\ref{lem:padding} with \(p=2/3\) and
\(p'=s_0=16/25=0.64\) to the learner in
Definition~\ref{def:streaming-model}.  This gives deterministic Borel
state-based rules and data-independent initialization, with uniform-prior
success at least \(s_0\); all remaining randomness is fixed.
Pre-generate all \(T\) rows, including rows after an early stop.

Put \(L_\epsilon=\log(1/\epsilon)\).  It suffices to consider dimensions
for which \(T\le dL_\epsilon\).  At other dimensions the claimed lower
bound already holds.  Lemma~\ref{lem:early-clock}, applied with
\(p_0=s_0\), gives
\begin{equation}\label{syn:output-count}
 s_0\le(T+1)2^M\epsilon^{d-1},\qquad
 L_\epsilon=o(d),\qquad T=o(d^2).
\end{equation}
Pad first to time \(T\) and then
to the least multiple of \(k\) at least \(T\), using the stopping-index
tags and ignored fresh pairs of Lemma~\ref{lem:padding}.  Its exact
per-index state count gives
\begin{equation}\label{syn:padded-width}
 K=(T+2)2^M,\qquad \log K=o(d^2).
\end{equation}

Since \(s_0>3/5\), Lemma~\ref{lem:raw-sample-bound} gives
\(T>\lfloor d/2\rfloor\) for all sufficiently large \(d\).  We use the
weaker consequence
\begin{equation}\label{syn:linear-endpoint}
 T>d/4.
\end{equation}

Let \(p=\lceil T/k\rceil\).  Write \(H_i\) for the tuple of states at the
ends of the first \(i\) blocks, with \(H_0\) constant.  Use one independent
\(B\sim\gamma_{m,d}\) and \(Y=BX\) throughout.  If a history
\(h\) has probability \(p_h>0\), Lemma~\ref{lem:conditional-priors}
gives its conditional signal density
\[
 \frac{\Prob\{H_{i-1}=h\mid X=x\}}{p_h}\le\frac1{p_h}
 \quad\text{relative to }\sigma_d.
\]
The next block matrix is independent of \((X,H_{i-1})\).  Given \(h\),
the initial state is fixed, and the next state is a measurable \(K\)-valued
function of the new block and its exact labels.  The side matrix remains
independent of this conditional experiment.  Proposition~\ref{syn:block}
therefore applies with \(\Lambda=1/p_h\).  Averaging it over histories gives
\begin{align}
 I(X;H_i\mid H_{i-1},B,Y)
 &\le\frac{\log K}{\ell}
   +C\left[d+\sum_h p_h\log(1+\log(1/p_h))\right] \notag\\
 &\le\frac{\log K}{\ell}
   +C[d+\log(1+(i-1)\log K)].
 \label{syn:history-step}
\end{align}
The last line uses concavity of \(t\mapsto\log(1+t)\) and
\(H(H_{i-1})\le(i-1)\log K\).  In particular it applies the bounded-prior
block theorem before averaging its density cost.  Since
\(p=o(d)\), \(\ell\ge d/16\), and \(\log K=o(d^2)\), the right side is
at most \(Cd\), with one absolute \(C\), for all sufficiently large \(d\).
The chain rule then gives
\begin{equation}\label{syn:history-information}
 I(X;H_p\mid B,Y)\le Cdp\le C(T+d).
\end{equation}
The postponed deterministic output is a function of the final tagged state,
and therefore of \(H_p\).

To obtain the matching endpoint, use Lemma~\ref{lem:residual-information}
with \(S=X\), \(V=H_p\), \(G=B\), and \(r_0=1/4\).
The output is a function of \(H_p\), and \(B\) is independent of the
entire signal--history experiment.  Moreover, \(2\epsilon<1/4\) and
\(m\le d/2\le d-2\).  For the residual radius
\(\rho_B=\norm{P_{\ker B}X}\), \eqref{eq:residual-markov} gives
\[
 \Prob\{\rho_B<1/4\}
 \le\frac{m/d}{1-1/16}\le\frac8{15}.
\]
Thus the lemma applies with success minus radius-failure probability at
least \(s_0-8/15=8/75\), and yields
\begin{align}
 I(X;H_p\mid B,Y)
 &\ge\frac8{75}(d-m-1)\log\frac1{8\epsilon}-\log2 \notag\\
 &\ge c_1dL_\epsilon
 \label{syn:history-lower}
\end{align}
for all sufficiently large \(d\), with an absolute \(c_1>0\).  Indeed,
\[
 \log\frac1{8\epsilon}
 \ge\left(1-\frac{\log8}{\log10}\right)L_\epsilon,
\]
and \(d-m-1\ge d/2-1\).

By \eqref{syn:linear-endpoint}, \(T+d<5T\).  Combining
\eqref{syn:history-information} and \eqref{syn:history-lower} therefore
gives \(T\ge c_2dL_\epsilon\) with an absolute \(c_2>0\).  The constants
are absolute; the eventual threshold comes only from \(M(d)=o(d^2)\).
All replicas and auxiliary projections have served solely to bound
information.  The learner's observations remain the prescribed fresh
Gaussian rows and exact labels.
\end{proof}
